\chapter{Klasyfikacja topologicznych rozmaitości czterowymiarowych}
\label{ch:freedman-klasyfikacja}
\index{Freedman!klasyfikacja wymiaru cztery}

W Rozdziale~\ref{ch:freedman-formy} obliczaliśmy formy
przecięcia i oglądaliśmy przeszkodę dla gładkiego
triku Whitneya. Uchwyty Cassona, opisane w
Rozdziale~\ref{ch:uchwyty-cassona}, stanowią drogę
do topologicznego twierdzenia o osadzaniu dysku.
Poniżej podamy precyzyjny wynik klasyfikacyjny
Freedmana jako wejście zewnętrzne i wyprowadzimy
z niego wnioski. Każde stwierdzenie o homeomorfizmie
dotyczy kategorii TOP; typy gładkie wymagają
osobnych argumentów.

\section{Niezmiennik Kirby'ego--Siebenmanna}
\label{sec:freedman-ks}
\index{niezmiennik Kirby'ego--Siebenmanna}

Niech $X$ będzie zamkniętą, spójną, topologiczną
rozmaitością $4$-wymiarową. Jej stabilna wiązka
styczna jest w kategorii TOP mikrowiązką. Teoria
Kirby'ego--Siebenmanna przypisuje jej klasę
\begin{equation}\label{eq:freedman-ks}
 \operatorname{ks}(X)\in H^4(X;\Z/2).
\end{equation}
Po ewaluacji na mod $2$ klasie fundamentalnej
otrzymujemy liczbę w $\Z/2$, również oznaczaną
$\operatorname{ks}(X)$. Intuicyjnie mierzy ona
przeszkodę do \emph{stabilnego} wygładzenia: dla
zamkniętej rozmaitości $4$-wymiarowej
$\operatorname{ks}(X)=0$ wtedy i tylko wtedy,
gdy $X\times\R$ dopuszcza strukturę gładką.
Jest to wynik zewnętrzny teorii wygładzania,
\href{https://archive.mpim-bonn.mpg.de/4789/2/FreedmanQuinn-TopologyOf4Manifolds-Reformatted2013.pdf}
{Freedman--Quinn, \S\,8.7--8.8 i \S\,10.2.2},
oparty na
\href{https://doi.org/10.1515/9781400881505}
{Kirby--Siebenmann, \emph{Foundational Essays on
Topological Manifolds, Smoothings, and Triangulations}
(1977)}.

W szczególności struktura gładka lub PL na $X$
wymusza $\operatorname{ks}(X)=0$. W wymiarze cztery
odwrotnej implikacji \emph{nie} wolno wyciągać z
tego twierdzenia: zerowanie przeszkody stabilnej
nie buduje automatycznie wygładzenia samego $X$.
Niezmiennik jest zachowywany przez homeomorfizmy,
a dla sum spójnych zachodzi
\[
 \operatorname{ks}(X\#Y)
   =\operatorname{ks}(X)+\operatorname{ks}(Y)
       \quad\text{w }\Z/2.
\]
Te własności również przyjmujemy z powyższego źródła.

Jeżeli $X$ jest dodatkowo jednospójna i zorientowana,
parzystość $Q_X$ odpowiada warunkowi
$w_2(X)=0$: jest to wzór Wu
$Q_X(x,x)\equiv\langle w_2(X),x\rangle\pmod2$
wraz z brakiem $2$-torsji w $H_1(X)$.
Jest to topologiczna wersja warunku spinowego;
również ją przyjmujemy jako wynik zewnętrzny,
\href{https://archive.mpim-bonn.mpg.de/4789/2/FreedmanQuinn-TopologyOf4Manifolds-Reformatted2013.pdf}
{Freedman--Quinn, \S\,10.2.1--10.2.2}.
Bez hipotezy o $H_1$ sama parzystość formy na
całkowitej homologii nie musi wykrywać całego $w_2$.

\section{Twierdzenie Freedmana jako wynik zewnętrzny}
\label{sec:freedman-classification}

\begin{twierdzenie}[Freedman: klasyfikacja jednospójna;
wynik zewnętrzny]
\label{thm:freedman-classification}
Niech $Q$ będzie symetryczną unimodularną formą
całkowitą na wolnej grupie abelowej skończonej rangi,
a $k\in\Z/2$. Jeśli $Q$ jest parzysta, załóżmy
dodatkowo
\begin{equation}\label{eq:freedman-even-ks}
 k\equiv\frac{\sigma(Q)}8\pmod2.
\end{equation}
Wtedy istnieje zamknięta, spójna, zorientowana,
jednospójna \emph{topologiczna} rozmaitość
$X^4$ z $Q_X\cong Q$ i $\operatorname{ks}(X)=k$.
Jeśli $Q$ jest nieparzysta, obie wartości $k=0,1$
są dopuszczalne.

Ponadto dwie takie rozmaitości są homeomorficzne
z zachowaniem orientacji wtedy i tylko wtedy,
gdy ich formy przecięcia są izometryczne i mają
ten sam niezmiennik Kirby'ego--Siebenmanna.
W wersji z ustaloną izometrią form homeomorfizm
można wybrać tak, by indukował właśnie tę izometrię.
\end{twierdzenie}

Wzór~\eqref{eq:freedman-even-ks} ma sens, ponieważ
sygnatura parzystej formy unimodularnej jest podzielna
przez $8$; jak zaznaczono po
Stwierdzeniu~\ref{prop:freedman-E8-properties},
ogólny fakt algebraiczny przyjmujemy z literatury.
Sam wzór i całe Twierdzenie~\ref{thm:freedman-classification}
przyjmujemy z
\href{https://projecteuclid.org/journals/journal-of-differential-geometry/volume-17/issue-3/The-topology-of-four-dimensional-manifolds/10.4310/jdg/1214437136.full}
{M.\ H.\ Freedman, \emph{The topology of four-dimensional
manifolds}, J. Differential Geom. 17 (1982), 357--453,
Twierdzenie 1.5}, w precyzyjnej postaci
\href{https://archive.mpim-bonn.mpg.de/4789/2/FreedmanQuinn-TopologyOf4Manifolds-Reformatted2013.pdf}
{M.\ H.\ Freedman i F.\ Quinn, \emph{Topology of
4-Manifolds}, \S\,10.1}.
W źródłach dowód korzysta z topologicznego osadzania
dysków i uchwytów Cassona. Poprzednie sekcje tego
rozdziału nie dowodzą istnienia ani jednoznaczności
homeomorfizmu z twierdzenia.

\begin{wniosek}[Topologiczna hipoteza Poincarégo w wymiarze cztery]
\label{cor:freedman-poincare-4}
Jeżeli zamknięta, spójna, jednospójna topologiczna
rozmaitość $X^4$ ma homologie całkowite sfery $S^4$,
to $X$ jest homeomorficzna z $S^4$.
W szczególności każda topologiczna sfera homotopijna
wymiaru cztery jest homeomorficzna z $S^4$.
\end{wniosek}
\begin{proof}
Warunek homologiczny daje $H_2(X;\Z)=0$, więc
$Q_X$ jest formą zerową na grupie rangi zero.
Jest parzysta, unimodularna w sensie pustej macierzy
o wyznaczniku $1$, i ma sygnaturę $0$.
Z~\eqref{eq:freedman-even-ks} wynika
$\operatorname{ks}(X)=0$. Standardowa $S^4$
ma te same dane, więc część jednoznaczności
Twierdzenia~\ref{thm:freedman-classification}
daje homeomorfizm. Wniosek nie stwierdza
istnienia dyfeomorfizmu.
\end{proof}

\begin{przyklad}[Cztery dane topologiczne]
\label{prz:freedman-examples}
Forma $H$ jest parzysta, ma sygnaturę $0$ i
$\operatorname{ks}=0$, więc jej jedynym zorientowanym
typem topologicznym jest $S^2\times S^2$.
Forma $[1]$ jest nieparzysta: istnieją dokładnie
dwa typy. Jeden to $\C P^2$ z
$\operatorname{ks}=0$; drugi ma $\operatorname{ks}=1$
i nie może mieć struktury gładkiej ani PL.

Forma $E_8$ jest parzysta i ma sygnaturę $8$.
Twierdzenie Freedmana daje dokładnie jeden
zorientowany typ topologiczny $X_{E_8}$ o tej formie,
a~\eqref{eq:freedman-even-ks} daje
$\operatorname{ks}(X_{E_8})=1$.
Wobec własności z Sekcji~\ref{sec:freedman-ks}
$X_{E_8}$ nie ma struktury gładkiej ani PL.
Natomiast forma $E_8\oplus E_8$ ma sygnaturę $16$,
więc odpowiadający jej typ topologiczny ma
$\operatorname{ks}=0$. Tego ostatniego rachunku
nie interpretujemy jako dowodu gładkości.
\end{przyklad}

Twierdzenie~\ref{thm:freedman-classification}
dotyczy tylko zamkniętych, zorientowanych,
jednospójnych rozmaitości w kategorii TOP.
Przypadek ogólnej grupy podstawowej wymaga
dodatkowych danych i hipotez. Własności gładkich
struktur na tej samej topologicznej rozmaitości
wymagają dalszych narzędzi, w szczególności
teorii Donaldsona; żaden rachunek samego
wyznacznika, parzystości i sygnatury ich nie rozstrzyga.

\par\smallskip
{\footnotesize\noindent\emph{Dalsza lektura:}
\href{https://projecteuclid.org/journals/journal-of-differential-geometry/volume-17/issue-3/The-topology-of-four-dimensional-manifolds/10.4310/jdg/1214437136.full}
{M.\ H.\ Freedman, \emph{The topology of four-dimensional
manifolds}, J. Differential Geom. 17 (1982), 357--453},
oryginalna konstrukcja;
\href{https://archive.mpim-bonn.mpg.de/4789/2/FreedmanQuinn-TopologyOf4Manifolds-Reformatted2013.pdf}
{M.\ H.\ Freedman i F.\ Quinn, \emph{Topology of
4-Manifolds} (1990)}, zwłaszcza \S\,10.1--10.2:
pełna postać klasyfikacji i niezmiennik
Kirby'ego--Siebenmanna.\par}
